\begin{frame}
	\fta{Einleitung}

	\s{
		Während die Berechnungen von Serienschaltungen ohne parallele geschaltete 
		Bestandteile recht aufwendungsarm zu bewerkstelligen sind, wird die Analyse von 
		Gleichstromnetzwerken mit Knoten, welche n Abzweigungen aufweisen, wobei die
		Anzahl der Zweige dabei $n \geq 2$ ist, bei steigender Knotenanzahl schnell 
		unübersichtlich. Eine vielschichtige Betrachtung des Gleichstromnetzwerkes ist 
		hier von Nöten. Für die Berechnung von Gleichstromnetzwerken stehen die 
		folgenden Methoden zur Verfügung: 
		}

	\b{
		In dem Modul 'Erweiterten Gleichstromnetzwerke' werden die folgenden Inhalte
		erläutert:
		}

	\begin{itemize}
		\item Knoten- und Maschenregel
		\item Superposition
		\item Knotenpotentialanalyse
		\item Maschenstromanalyse
	\end{itemize}

	\s{
		Auch die Analyse von elektrischen Netzwerken mit mehr als einer Strom- 
		oder Spannungsquelle ist nicht trivial. Die Bestimmung aller Ströme und 
		Spannungen des elektrischen Netzwerkes aus Abbildung \ref{BildEinleitungsnetzwerk} 
		ist allein durch die Berechnungen von Serienschaltungen und Parallelschaltungen nicht 
		möglich. 
		}
 
	\fu{
		\resizebox{.5\textwidth}{!}{%
		\input{Tikz/src/beispielnetzwerk.tex}\alttext{Die Abbildung zeigt ein verzweigtes Stromnetzwerk, dessen Leitungen einen annähernd ovalen Rahmen bilden, der oben und unten jeweils über ein waagerechtes Bauteil geschlossen ist. Entlang dieses Rahmens sind mehrere rechteckige Symbole (z. B. Widerstände) sowohl in Reihe als auch in schrägen Querzweigen angeordnet, sodass mehrere Parallelpfade entstehen. In das Netz sind vier kreisförmige Symbole (z. B. Spannungsquellen oder Messgeräte) eingebunden: etwa auf der linken Seite, unten nahe der Mitte, oben nahe der Mitte und rechts unten; an diesen Punkten geben kleine Pfeile die jeweilige Stromrichtung oder Spannung nach oben oder unten an.}%
		}
	}{{\bf Beispielnetzwerk.} Elektrisches Gleichstromnetzwerk mit einer Vielzahl an Widerständen,
	Stromquellen und Spannungsquellen. \label{BildEinleitungsnetzwerk}}
\end{frame}